\section{Introduction}
\label{sec:intro}

Quarterly 13F filings are the most complete public record of
institutional equity positioning, and the most heavily mined. The
literature's box score is mixed by construction: the filings are
stale (up to 45 days), long-only, and quarterly, so most of what they
reveal is already in prices. Our starting premise is therefore not
``institutions know something,'' but a narrower mechanical claim:
\emph{the filings occasionally reveal actions that were either costly
to take or impossible to avoid}, and only those actions should carry
a premium. A manager displacing an existing top-decile position to
make room for a new name has paid a real price for that opinion; a
manager liquidating a third of a book under redemptions is not
expressing an opinion at all. Everything in this report is organized
around whether a signal isolates one of those two situations.

The work divides into three parts. The first
(Sections~\ref{sec:data}--\ref{sec:universe}) is data engineering,
and we give it the space usually reserved for results because in this
dataset the engineering \emph{is} the result: a bitemporal
point-in-time store keyed on filing acceptance timestamps; amendment
semantics read from cover pages rather than guessed; per-filing value
scale detection; a CUSIP$\to$issuer-name$\to$current-ticker crosswalk
with human resolution of edge cases; and a value-weighted coverage
audit that late in the project found 46.4\% of book value unmapped
--- EDGAR's name abbreviations defeat exact matching precisely on the
mega-caps --- forcing a fix and a 24-run re-verification campaign in
which one previously reported result was retracted
(Section~\ref{sec:campaign}). We report the retraction with its
mechanism because the episode is the strongest evidence we can offer
that the rest of the numbers mean something.

The second part (Sections~\ref{sec:taxonomy}--\ref{sec:fss}) is a
taxonomy: six classes of positioning signals, seventeen catalogue
constructions evaluated under one fixed protocol, five pre-registered
interaction effects, and the two survivors. New-Conviction Breadth
(NCB) counts, per stock, the fraction of filers that opened or grew
($\ge 25\%$, split-adjusted) a position sitting in the top decile of
their own book, residualized on size and dollar liquidity because a
formal placebo (``dart-throwing value-weighted managers'') shows the
raw count is mechanically correlated with both. Forced-Sale Supply
(FSS) converts a validated manager-level flow instrument into a
stock-level supply measure. For each we present the economic
mechanism first, the returns second, and the hyperparameter surface
third --- every parameter in this report is either justified by a
measurement or explicitly labeled a convention.

The third part (Sections~\ref{sec:models}--\ref{sec:limitations})
assembles the pieces into production objects and stress-tests the
assembly: a Fama--MacBeth expected-return combiner over ten
characteristics (the production model), an IPCA challenger under the
\citet{kps2019} instrumented-loadings framework, a gradient-boosted
ex-post exercise that rediscovers NCB and finds no nonlinear premium,
net-of-cost engine results with explicit commission and borrow
accounting, and the descriptive context --- the median large 13F
filer underperforms the S\&P~500 by $3.3\%$ per year before fees ---
that explains why imitating these books is a losing design and
harvesting their \emph{flows} is not.

Two regularities recur so often that we state them as laws and tag
every result with them. \textbf{Law 1: states have no arrow.} Levels
of ownership, breadth, crowding, concentration, or factor tilt never
predict returns in this data, at any clock, in fifteen independent
formulations; changes and constrained actions do. \textbf{Law 2:
distress is information when observed, never when predicted.}
Contemporaneous forced selling is measurable with $t\approx+10$;
every attempt to anticipate it --- daily synthetic stress, quarterly
flow forecasting, single-name crash cascades --- fails its mechanism
gate or dies on the full panel.
